% ================= 讲义导言区（与主文件配套，xelatex + ctex） =================
% 主文件写法：
%   \documentclass[11pt,a4paper]{article}
%   % ================= 讲义导言区（与主文件配套，xelatex + ctex） =================
% 主文件写法：
%   \documentclass[11pt,a4paper]{article}
%   % ================= 讲义导言区（与主文件配套，xelatex + ctex） =================
% 主文件写法：
%   \documentclass[11pt,a4paper]{article}
%   % ================= 讲义导言区（与主文件配套，xelatex + ctex） =================
% 主文件写法：
%   \documentclass[11pt,a4paper]{article}
%   \input{preamble.tex}   ← 本文件（只含 packages 与格式定义）
%   \title{...}
%   \begin{document} ... \tableofcontents ... 各章 ... \end{document}
% 注意：xelatex 必须在讲义目录内运行（不要用 -output-directory，会破坏 figs/ 相对路径）

\usepackage[UTF8]{ctex}
\usepackage{amsmath,amssymb,amsthm}
\usepackage{graphicx}
\usepackage{multirow}
\usepackage{array}
\usepackage{booktabs}
\usepackage{longtable}
\usepackage{enumitem}
\usepackage{geometry}
\usepackage{caption}
\usepackage{float}
\usepackage{fancyhdr}
\usepackage{xcolor}
\usepackage{titlesec}
\usepackage{tcolorbox}
\tcbuselibrary{breakable,skins}   % ← breakable 必须显式加载，否则报错
\usepackage{makecell}
\usepackage{hyperref}

\geometry{a4paper,left=2.3cm,right=2.3cm,top=2.5cm,bottom=2.3cm}

% ---- 颜色 ----
\definecolor{mainblue}{RGB}{20,60,130}
\definecolor{deepblue}{RGB}{0,55,120}
\definecolor{redbold}{RGB}{180,20,20}
\definecolor{boxbg}{RGB}{242,246,252}
\definecolor{boxframe}{RGB}{60,110,180}
\definecolor{thmframe}{RGB}{160,60,40}
\definecolor{thmbg}{RGB}{253,247,244}

\hypersetup{colorlinks=true,linkcolor=deepblue,urlcolor=deepblue,citecolor=deepblue,
  bookmarksnumbered=true,bookmarksopen=true}

% ---- 标题格式 ----
\titleformat{\section}{\normalfont\Large\bfseries\color{mainblue}}{\thesection}{0.6em}{}
\titleformat{\subsection}{\normalfont\large\bfseries\color{deepblue}}{\thesubsection}{0.5em}{}
\titleformat{\subsubsection}{\normalfont\normalsize\bfseries\color{boxframe}}{\thesubsubsection}{0.5em}{}
\setcounter{tocdepth}{2}
\setcounter{secnumdepth}{3}

% ---- 页眉页脚 ----
\pagestyle{fancy}
\fancyhf{}
\fancyhead[L]{\small\color{gray}课程讲义}
\fancyhead[R]{\small\color{gray}\leftmark}
\fancyfoot[C]{\small\thepage}
\renewcommand{\headrulewidth}{0.4pt}

% ---- 彩色提示框 ----
\newtcolorbox{keybox}[1][]{
  colback=boxbg, colframe=boxframe, boxrule=0.8pt, arc=2pt,
  left=6pt,right=6pt,top=5pt,bottom=5pt,
  fonttitle=\bfseries\color{white}, coltitle=white,
  title=#1, breakable}

\newtcolorbox{exbox}[1][]{
  colback=thmbg, colframe=thmframe, boxrule=0.8pt, arc=2pt,
  left=6pt,right=6pt,top=5pt,bottom=5pt,
  fonttitle=\bfseries\color{white}, coltitle=white,
  title=#1, breakable}

% ---- 例题 / 习题环境 ----
\newenvironment{example}[1]{\par\medskip\noindent\textbf{\color{redbold}【例】#1}\par\smallskip}{\par\medskip}
\newenvironment{solution}{\par\noindent\textbf{\color{deepblue}【解】}\par\smallskip}{\par\medskip}
\newenvironment{exercise}[1]{\par\medskip\noindent\textbf{\color{mainblue}【习题 #1】}\par\smallskip}{\par\medskip}
\newenvironment{exsolution}{\par\noindent\textbf{\color{deepblue}【解答】}\par\smallskip}{\par\medskip}

% ---- 常用宏 ----
\newcommand{\ov}[1]{\overline{#1}}
\newcommand{\XOR}{\oplus}
\newcommand{\term}[1]{\textbf{\color{redbold}#1}}

% ---- 图表 ----
\captionsetup{font=small,labelfont=bf,skip=4pt}
\renewcommand{\arraystretch}{1.15}
% 电路图统一样式：按原比例，限制最大宽度与高度（防止大图溢出页面）
\newcommand{\circfig}[3][1.0]{%
  \begin{center}
  \includegraphics[width=#1\linewidth,height=0.42\textheight,keepaspectratio]{#2}
  \captionof{figure}{#3}
  \end{center}}
   ← 本文件（只含 packages 与格式定义）
%   \title{...}
%   \begin{document} ... \tableofcontents ... 各章 ... \end{document}
% 注意：xelatex 必须在讲义目录内运行（不要用 -output-directory，会破坏 figs/ 相对路径）

\usepackage[UTF8]{ctex}
\usepackage{amsmath,amssymb,amsthm}
\usepackage{graphicx}
\usepackage{multirow}
\usepackage{array}
\usepackage{booktabs}
\usepackage{longtable}
\usepackage{enumitem}
\usepackage{geometry}
\usepackage{caption}
\usepackage{float}
\usepackage{fancyhdr}
\usepackage{xcolor}
\usepackage{titlesec}
\usepackage{tcolorbox}
\tcbuselibrary{breakable,skins}   % ← breakable 必须显式加载，否则报错
\usepackage{makecell}
\usepackage{hyperref}

\geometry{a4paper,left=2.3cm,right=2.3cm,top=2.5cm,bottom=2.3cm}

% ---- 颜色 ----
\definecolor{mainblue}{RGB}{20,60,130}
\definecolor{deepblue}{RGB}{0,55,120}
\definecolor{redbold}{RGB}{180,20,20}
\definecolor{boxbg}{RGB}{242,246,252}
\definecolor{boxframe}{RGB}{60,110,180}
\definecolor{thmframe}{RGB}{160,60,40}
\definecolor{thmbg}{RGB}{253,247,244}

\hypersetup{colorlinks=true,linkcolor=deepblue,urlcolor=deepblue,citecolor=deepblue,
  bookmarksnumbered=true,bookmarksopen=true}

% ---- 标题格式 ----
\titleformat{\section}{\normalfont\Large\bfseries\color{mainblue}}{\thesection}{0.6em}{}
\titleformat{\subsection}{\normalfont\large\bfseries\color{deepblue}}{\thesubsection}{0.5em}{}
\titleformat{\subsubsection}{\normalfont\normalsize\bfseries\color{boxframe}}{\thesubsubsection}{0.5em}{}
\setcounter{tocdepth}{2}
\setcounter{secnumdepth}{3}

% ---- 页眉页脚 ----
\pagestyle{fancy}
\fancyhf{}
\fancyhead[L]{\small\color{gray}课程讲义}
\fancyhead[R]{\small\color{gray}\leftmark}
\fancyfoot[C]{\small\thepage}
\renewcommand{\headrulewidth}{0.4pt}

% ---- 彩色提示框 ----
\newtcolorbox{keybox}[1][]{
  colback=boxbg, colframe=boxframe, boxrule=0.8pt, arc=2pt,
  left=6pt,right=6pt,top=5pt,bottom=5pt,
  fonttitle=\bfseries\color{white}, coltitle=white,
  title=#1, breakable}

\newtcolorbox{exbox}[1][]{
  colback=thmbg, colframe=thmframe, boxrule=0.8pt, arc=2pt,
  left=6pt,right=6pt,top=5pt,bottom=5pt,
  fonttitle=\bfseries\color{white}, coltitle=white,
  title=#1, breakable}

% ---- 例题 / 习题环境 ----
\newenvironment{example}[1]{\par\medskip\noindent\textbf{\color{redbold}【例】#1}\par\smallskip}{\par\medskip}
\newenvironment{solution}{\par\noindent\textbf{\color{deepblue}【解】}\par\smallskip}{\par\medskip}
\newenvironment{exercise}[1]{\par\medskip\noindent\textbf{\color{mainblue}【习题 #1】}\par\smallskip}{\par\medskip}
\newenvironment{exsolution}{\par\noindent\textbf{\color{deepblue}【解答】}\par\smallskip}{\par\medskip}

% ---- 常用宏 ----
\newcommand{\ov}[1]{\overline{#1}}
\newcommand{\XOR}{\oplus}
\newcommand{\term}[1]{\textbf{\color{redbold}#1}}

% ---- 图表 ----
\captionsetup{font=small,labelfont=bf,skip=4pt}
\renewcommand{\arraystretch}{1.15}
% 电路图统一样式：按原比例，限制最大宽度与高度（防止大图溢出页面）
\newcommand{\circfig}[3][1.0]{%
  \begin{center}
  \includegraphics[width=#1\linewidth,height=0.42\textheight,keepaspectratio]{#2}
  \captionof{figure}{#3}
  \end{center}}
   ← 本文件（只含 packages 与格式定义）
%   \title{...}
%   \begin{document} ... \tableofcontents ... 各章 ... \end{document}
% 注意：xelatex 必须在讲义目录内运行（不要用 -output-directory，会破坏 figs/ 相对路径）

\usepackage[UTF8]{ctex}
\usepackage{amsmath,amssymb,amsthm}
\usepackage{graphicx}
\usepackage{multirow}
\usepackage{array}
\usepackage{booktabs}
\usepackage{longtable}
\usepackage{enumitem}
\usepackage{geometry}
\usepackage{caption}
\usepackage{float}
\usepackage{fancyhdr}
\usepackage{xcolor}
\usepackage{titlesec}
\usepackage{tcolorbox}
\tcbuselibrary{breakable,skins}   % ← breakable 必须显式加载，否则报错
\usepackage{makecell}
\usepackage{hyperref}

\geometry{a4paper,left=2.3cm,right=2.3cm,top=2.5cm,bottom=2.3cm}

% ---- 颜色 ----
\definecolor{mainblue}{RGB}{20,60,130}
\definecolor{deepblue}{RGB}{0,55,120}
\definecolor{redbold}{RGB}{180,20,20}
\definecolor{boxbg}{RGB}{242,246,252}
\definecolor{boxframe}{RGB}{60,110,180}
\definecolor{thmframe}{RGB}{160,60,40}
\definecolor{thmbg}{RGB}{253,247,244}

\hypersetup{colorlinks=true,linkcolor=deepblue,urlcolor=deepblue,citecolor=deepblue,
  bookmarksnumbered=true,bookmarksopen=true}

% ---- 标题格式 ----
\titleformat{\section}{\normalfont\Large\bfseries\color{mainblue}}{\thesection}{0.6em}{}
\titleformat{\subsection}{\normalfont\large\bfseries\color{deepblue}}{\thesubsection}{0.5em}{}
\titleformat{\subsubsection}{\normalfont\normalsize\bfseries\color{boxframe}}{\thesubsubsection}{0.5em}{}
\setcounter{tocdepth}{2}
\setcounter{secnumdepth}{3}

% ---- 页眉页脚 ----
\pagestyle{fancy}
\fancyhf{}
\fancyhead[L]{\small\color{gray}课程讲义}
\fancyhead[R]{\small\color{gray}\leftmark}
\fancyfoot[C]{\small\thepage}
\renewcommand{\headrulewidth}{0.4pt}

% ---- 彩色提示框 ----
\newtcolorbox{keybox}[1][]{
  colback=boxbg, colframe=boxframe, boxrule=0.8pt, arc=2pt,
  left=6pt,right=6pt,top=5pt,bottom=5pt,
  fonttitle=\bfseries\color{white}, coltitle=white,
  title=#1, breakable}

\newtcolorbox{exbox}[1][]{
  colback=thmbg, colframe=thmframe, boxrule=0.8pt, arc=2pt,
  left=6pt,right=6pt,top=5pt,bottom=5pt,
  fonttitle=\bfseries\color{white}, coltitle=white,
  title=#1, breakable}

% ---- 例题 / 习题环境 ----
\newenvironment{example}[1]{\par\medskip\noindent\textbf{\color{redbold}【例】#1}\par\smallskip}{\par\medskip}
\newenvironment{solution}{\par\noindent\textbf{\color{deepblue}【解】}\par\smallskip}{\par\medskip}
\newenvironment{exercise}[1]{\par\medskip\noindent\textbf{\color{mainblue}【习题 #1】}\par\smallskip}{\par\medskip}
\newenvironment{exsolution}{\par\noindent\textbf{\color{deepblue}【解答】}\par\smallskip}{\par\medskip}

% ---- 常用宏 ----
\newcommand{\ov}[1]{\overline{#1}}
\newcommand{\XOR}{\oplus}
\newcommand{\term}[1]{\textbf{\color{redbold}#1}}

% ---- 图表 ----
\captionsetup{font=small,labelfont=bf,skip=4pt}
\renewcommand{\arraystretch}{1.15}
% 电路图统一样式：按原比例，限制最大宽度与高度（防止大图溢出页面）
\newcommand{\circfig}[3][1.0]{%
  \begin{center}
  \includegraphics[width=#1\linewidth,height=0.42\textheight,keepaspectratio]{#2}
  \captionof{figure}{#3}
  \end{center}}
   ← 本文件（只含 packages 与格式定义）
%   \title{...}
%   \begin{document} ... \tableofcontents ... 各章 ... \end{document}
% 注意：xelatex 必须在讲义目录内运行（不要用 -output-directory，会破坏 figs/ 相对路径）

\usepackage[UTF8]{ctex}
\usepackage{amsmath,amssymb,amsthm}
\usepackage{graphicx}
\usepackage{multirow}
\usepackage{array}
\usepackage{booktabs}
\usepackage{longtable}
\usepackage{enumitem}
\usepackage{geometry}
\usepackage{caption}
\usepackage{float}
\usepackage{fancyhdr}
\usepackage{xcolor}
\usepackage{titlesec}
\usepackage{tcolorbox}
\tcbuselibrary{breakable,skins}   % ← breakable 必须显式加载，否则报错
\usepackage{makecell}
\usepackage{hyperref}

\geometry{a4paper,left=2.3cm,right=2.3cm,top=2.5cm,bottom=2.3cm}

% ---- 颜色 ----
\definecolor{mainblue}{RGB}{20,60,130}
\definecolor{deepblue}{RGB}{0,55,120}
\definecolor{redbold}{RGB}{180,20,20}
\definecolor{boxbg}{RGB}{242,246,252}
\definecolor{boxframe}{RGB}{60,110,180}
\definecolor{thmframe}{RGB}{160,60,40}
\definecolor{thmbg}{RGB}{253,247,244}

\hypersetup{colorlinks=true,linkcolor=deepblue,urlcolor=deepblue,citecolor=deepblue,
  bookmarksnumbered=true,bookmarksopen=true}

% ---- 标题格式 ----
\titleformat{\section}{\normalfont\Large\bfseries\color{mainblue}}{\thesection}{0.6em}{}
\titleformat{\subsection}{\normalfont\large\bfseries\color{deepblue}}{\thesubsection}{0.5em}{}
\titleformat{\subsubsection}{\normalfont\normalsize\bfseries\color{boxframe}}{\thesubsubsection}{0.5em}{}
\setcounter{tocdepth}{2}
\setcounter{secnumdepth}{3}

% ---- 页眉页脚 ----
\pagestyle{fancy}
\fancyhf{}
\fancyhead[L]{\small\color{gray}课程讲义}
\fancyhead[R]{\small\color{gray}\leftmark}
\fancyfoot[C]{\small\thepage}
\renewcommand{\headrulewidth}{0.4pt}

% ---- 彩色提示框 ----
\newtcolorbox{keybox}[1][]{
  colback=boxbg, colframe=boxframe, boxrule=0.8pt, arc=2pt,
  left=6pt,right=6pt,top=5pt,bottom=5pt,
  fonttitle=\bfseries\color{white}, coltitle=white,
  title=#1, breakable}

\newtcolorbox{exbox}[1][]{
  colback=thmbg, colframe=thmframe, boxrule=0.8pt, arc=2pt,
  left=6pt,right=6pt,top=5pt,bottom=5pt,
  fonttitle=\bfseries\color{white}, coltitle=white,
  title=#1, breakable}

% ---- 例题 / 习题环境 ----
\newenvironment{example}[1]{\par\medskip\noindent\textbf{\color{redbold}【例】#1}\par\smallskip}{\par\medskip}
\newenvironment{solution}{\par\noindent\textbf{\color{deepblue}【解】}\par\smallskip}{\par\medskip}
\newenvironment{exercise}[1]{\par\medskip\noindent\textbf{\color{mainblue}【习题 #1】}\par\smallskip}{\par\medskip}
\newenvironment{exsolution}{\par\noindent\textbf{\color{deepblue}【解答】}\par\smallskip}{\par\medskip}

% ---- 常用宏 ----
\newcommand{\ov}[1]{\overline{#1}}
\newcommand{\XOR}{\oplus}
\newcommand{\term}[1]{\textbf{\color{redbold}#1}}

% ---- 图表 ----
\captionsetup{font=small,labelfont=bf,skip=4pt}
\renewcommand{\arraystretch}{1.15}
% 电路图统一样式：按原比例，限制最大宽度与高度（防止大图溢出页面）
\newcommand{\circfig}[3][1.0]{%
  \begin{center}
  \includegraphics[width=#1\linewidth,height=0.42\textheight,keepaspectratio]{#2}
  \captionof{figure}{#3}
  \end{center}}
